\begin{table}[H]
\caption{Mô tả Use Case: Xem danh sách nội dung}
\label{usecase-3-01}
\centering
\renewcommand{\arraystretch}{1.2}
\begin{tabular}{@{}p{4cm} p{11cm}@{}}
\toprule
\textbf{Use case ID}       & UC-M3-01 \\ \midrule
\textbf{Tên Use case}      & Xem danh sách nội dung \\ \midrule
\textbf{Các tác nhân}      & Traveler, Guest, Provider \\ \midrule
\textbf{Mô tả}             & Cho phép người dùng xem danh sách bài viết hiển thị trên trang chủ hoặc bảng tin khám phá hoặc theo chủ đề. \\ \midrule
\textbf{Tiền điều kiện}    & Hệ thống đã có sẵn dữ liệu nội dung. \\ \midrule
\textbf{Hậu điều kiện}     & Danh sách nội dung được hiển thị theo bộ lọc mặc định hoặc tùy chọn người dùng. \\ \midrule
\textbf{Luồng sự kiện chính} & 
\begin{tabular}[t]{@{}p{10.5cm}@{}}
1. Người dùng truy cập trang chủ hoặc bảng tin khám phá.\\
2. Hệ thống hiển thị danh sách nội dung công khai theo thứ tự mặc định.
\end{tabular} \\ \midrule
\textbf{Luồng ngoại lệ} & 2a. Nếu hệ thống bị mất kết nối, màn hình hiển thị “Lỗi kết nối” \\ 
\bottomrule
\end{tabular}
\end{table}

\begin{table}[H]
\caption{Mô tả Use Case: Lọc và sắp xếp nội dung}
\label{usecase-3-02}
\centering
\renewcommand{\arraystretch}{1.2}
\begin{tabular}{@{}p{4cm} p{11cm}@{}}
\toprule
\textbf{Use case ID}       & UC-M3-02 \\ \midrule
\textbf{Tên Use case}      & Lọc và sắp xếp nội dung \\ \midrule
\textbf{Các tác nhân}      & Traveler, Guest, Provider \\ \midrule
\textbf{Mô tả}             & Người dùng sắp xếp danh sách theo tiêu chí như: độ phổ biến, đánh giá cao, hoặc lọc danh sách theo các tiêu chí như: xếp hạng theo đánh giá khách hàng, khu vực, các bộ lọc đặc biệt (phù hợp với gia đình, cho phép thú cưng, phù hợp với solo-trip,...) \\ \midrule
\textbf{Trigger}    & Người dùng chọn mục "Sắp xếp theo…", "Lọc nâng cao". \\ \midrule
\textbf{Tiền điều kiện}    & Hệ thống đã có sẵn dữ liệu nội dung. \\ \midrule
\textbf{Hậu điều kiện}     & Hệ thống hiển thị danh sách kết quả tìm kiếm theo phân trang. \\ \midrule
\textbf{Luồng sự kiện chính} & 
\begin{tabular}[t]{@{}p{10.5cm}@{}}
1. Người chọn các tiêu chí lọc, sắp xếp.\\
2. Người dùng nhấn nút "Áp dụng".\\
3. Hệ thống truy vấn dữ liệu và hiển thị kết quả.
\end{tabular} \\ \midrule
\textbf{Luồng ngoại lệ} & \begin{tabular}[t]{@{}p{10.5cm}@{}}
3a. Không tìm thấy kết quả → hệ thống hiển thị thông báo "Không tìm thấy nội dung phù hợp."\\
3b. Lỗi kết nối → hiển thị "Không thể tải dữ liệu, vui lòng thử lại."
\end{tabular} \\ 
\bottomrule
\end{tabular}
\end{table}


\begin{table}[H]
\caption{Mô tả Use Case: Tìm kiếm nội dung}
\label{usecase-3-04}
\centering
\renewcommand{\arraystretch}{1.2}
\begin{tabular}{@{}p{4cm} p{11cm}@{}}
\toprule
\textbf{Use case ID}       & UC-M3-04 \\ \midrule
\textbf{Tên Use case}      & Tìm kiếm nội dung \\ \midrule
\textbf{Các tác nhân}      & Traveler, Guest, Provider \\ \midrule
\textbf{Mô tả}             & Cho phép người dùng tìm kiếm bài viết, địa điểm hoặc chủ đề theo từ khóa, danh mục. \\ \midrule
\textbf{Tiền điều kiện}    & Người dùng có thể truy cập thanh tìm kiếm. \\ \midrule
\textbf{Hậu điều kiện}     & Hệ thống hiển thị danh sách kết quả tìm kiếm theo phân trang. \\ \midrule
\textbf{Luồng sự kiện chính} & 
\begin{tabular}[t]{@{}p{10.5cm}@{}}
1. Người nhập từ khóa tìm kiếm \\
2. Người dùng nhấn nút tìm kiếm \\
3. Hệ thống truy vấn dữ liệu và hiển thị kết quả.
\end{tabular} \\ \midrule
\textbf{Luồng ngoại lệ} &
\begin{tabular}[t]{@{}p{10.5cm}@{}}
3a. Không tìm thấy kết quả → hệ thống hiển thị thông báo "Không tìm thấy nội dung phù hợp với từ khóa." \\
3b. Lỗi kết nối → hiển thị "Không thể tải dữ liệu, vui lòng thử lại."
\end{tabular} \\ \midrule
\textbf{Luồng thay thế} &
\begin{tabular}[t]{@{}p{10.5cm}@{}}
1. Người chọn danh mục/chủ đề gợi ý trên thanh công cụ tìm kiếm. \\
2. Người dùng nhấn nút tìm kiếm. \\
3. Hệ thống xử lý và hiển thị kết quả. 
\end{tabular} \\ 
\bottomrule
\end{tabular}
\end{table}

\begin{table}[H]
\caption{Mô tả Use Case: Lưu địa điểm, bài viết}
\label{usecase-3-05}
\centering
\renewcommand{\arraystretch}{1.2}
\begin{tabular}{@{}p{4cm} p{11cm}@{}}
\toprule
\textbf{Use case ID}       & UC-M3-05 \\ \midrule
\textbf{Tên Use case}      & Lưu địa điểm, bài viết \\ \midrule
\textbf{Các tác nhân}      & Traveler \\ \midrule
\textbf{Mô tả}             & Cho phép người dùng (Traveler) lưu lại bài viết hoặc địa điểm yêu thích để xem lại sau. \\ \midrule
\textbf{Trigger}    & Người dùng nhấn biểu tượng “Lưu” trên bài viết hoặc địa điểm. \\ \midrule
\textbf{Tiền điều kiện}    & Người dùng đã đăng nhập. \\ \midrule
\textbf{Hậu điều kiện}     & Nội dung được thêm vào danh sách đã lưu. \\ \midrule
\textbf{Luồng sự kiện chính} & 
\begin{tabular}[t]{@{}p{10.5cm}@{}}
1. Người dùng chọn biểu tượng "Lưu" \space tại bài viết hoặc địa điểm.\\
2. Hệ thống hiển thị danh sách "Bộ sưu tập" \space của người dùng.\\
3. Người dùng tuỳ chọn "Bộ sưu tập" \space để lưu bài viết/ địa điểm.\\
4. Hệ thống ghi nhận và thêm nội dung vào danh sách đã lưu.
\end{tabular} \\ \midrule
\textbf{Luồng ngoại lệ} &
\begin{tabular}[t]{@{}p{10.5cm}@{}}
1a. Chưa đăng nhập → hiển thị "Vui lòng đăng nhập để sử dụng tính năng lưu."
\end{tabular} \\ \midrule
\textbf{Luồng thay thế} &
\begin{tabular}[t]{@{}p{10.5cm}@{}}
3(a)1.  Người dùng chọn biểu tượng "+" \space để tạo "Bộ sưu tập".\\  
2. Hệ thống tự động lưu bài viết/ địa điểm vào bộ sưu tập vừa tạo.\\
3(b)1. Hệ thống tự động lưu bài viết/địa điểm vào danh mục mặc định "Tất cả bài viết".
\end{tabular} \\ 
\bottomrule
\end{tabular}
\end{table}

\begin{table}[H]
\caption{Mô tả Use Case: Xem danh sách nội dung địa điểm, bài viết đã lưu}
\label{usecase-3-06}
\centering
\renewcommand{\arraystretch}{1.2}
\begin{tabular}{@{}p{4cm} p{11cm}@{}}
\toprule
\textbf{Use case ID}       & UC-M3-06 \\ \midrule
\textbf{Tên Use case}      & Xem danh sách nội dung địa điểm, bài viết đã lưu \\ \midrule
\textbf{Các tác nhân}      & Traveler \\ \midrule
\textbf{Mô tả}             & Cho phép người dùng (Traveler) xem lại bài viết hoặc địa điểm yêu thích.\\ \midrule
\textbf{Trigger}    & Người dùng chọn biểu tượng “Bộ sưu tập” trên trang cá nhân \\ \midrule
\textbf{Tiền điều kiện}    & 
\begin{tabular}[t]{@{}p{10.5cm}@{}}
1. Người dùng đã đăng nhập.\\
2. Người dùng đã lưu ít nhất một nội dung.
\end{tabular} \\ \midrule
\textbf{Hậu điều kiện}     & Danh sách nội dung được hiển thị theo thứ tự mặc định (mới nhất -> cũ nhất) \\ \midrule
\textbf{Luồng sự kiện chính} & 
\begin{tabular}[t]{@{}p{10.5cm}@{}}
1. Người dùng chọn biểu tượng "Bộ sưu tập".\\
2. Hệ thống hiển thị các nội dung đã lưu theo các bộ sưu tập.\\
3. Người dùng chọn xem "tất cả bài viết" \space hoặc "bộ sưu tập" \space cụ thể.\\
4. Hệ thống hiển thị danh sách chi tiết.
\end{tabular} \\ \midrule
\textbf{Luồng ngoại lệ} &
\begin{tabular}[t]{@{}p{10.5cm}@{}}
2a. Danh sách trống → hiển thị "Bạn chưa lưu nội dung nào."
\end{tabular} \\
\bottomrule
\end{tabular}
\end{table}

\begin{table}[H]
\caption{Mô tả Use Case: Theo dõi chủ đề hoặc địa điểm}
\label{usecase-3-10}
\centering
\renewcommand{\arraystretch}{1.2}
\begin{tabular}{@{}p{4cm} p{11cm}@{}}
\toprule
\textbf{Use case ID}       & UC-M3-10 \\ \midrule
\textbf{Tên Use case}      & Theo dõi chủ đề hoặc địa điểm \\ \midrule
\textbf{Các tác nhân}      & Traveler \\ \midrule
\textbf{Mô tả}             & Cho phép người dùng theo dõi hoặc hủy theo dõi một chủ đề hoặc địa điểm.\\ \midrule
\textbf{Trigger}    & Người dùng nhấn nút "Theo dõi".\\ \midrule
\textbf{Tiền điều kiện}    & Người dùng đã đăng nhập.\\ \midrule
\textbf{Hậu điều kiện}     & Danh sách theo dõi được cập nhật \\ \midrule
\textbf{Luồng sự kiện chính} & 
\begin{tabular}[t]{@{}p{10.5cm}@{}}
1. Người dùng chọn nút "Theo dõi" \space chủ đề hoặc địa điểm. \\
2. Hệ thống ghi nhận hành động và cập nhật trạng thái theo dõi.
\end{tabular} \\ \midrule
\textbf{Luồng ngoại lệ} &
\begin{tabular}[t]{@{}p{10.5cm}@{}}
1a. Người dùng chưa đăng nhập → hiển thị "Vui lòng đăng nhập để theo dõi nội dung."
\end{tabular} \\
\bottomrule
\end{tabular}
\end{table}
